\chapter{Testing}
\label{cap:testing}

\begin{inbreve}
Il testing esegue un programma su dati di prova e controlla che si comporti come previsto. Il capitolo definisce i
tipi di testing e segue il \emph{functional testing} nelle sue fasi: unit, feature, system e release testing. Poi
spiega come automatizzare i test, il test-driven development e il security testing. Si chiude con i limiti del testing
e con le code review, che lo completano. Presuppone i capitoli~\ref{cap:agile} e~\ref{cap:feature}.
\end{inbreve}

\section{Che cos'è il testing}
\label{sec:test-intro}

\begin{definizione}[Testing]
Il \termine{testing} consiste in due passi:
\begin{enumerate}
  \item eseguire il programma con dati che simulano gli input degli utenti;
  \item osservare il comportamento del programma.
\end{enumerate}
Se il comportamento è quello atteso il test \emph{riesce}, altrimenti \emph{fallisce} (\figrif{fig:test-ciclo}).
\end{definizione}

\begin{figure}[ht]
\centering
\begin{tikzpicture}
\node[blocco, text width=3.4cm] (e) at (0,0) {1. Eseguire il programma con dati che simulano gli input};
\node[blocco, text width=2.6cm] (o) at (4.4,0) {2. Osservare il comportamento};
\node[draw, diamond, aspect=2.2, inner sep=1pt, font=\small, align=center] (q) at (8.6,0) {comportamento\\atteso?};
\node[blocco] (ok) at (8.6,1.8) {test riuscito};
\node[blocco] (ko) at (8.6,-1.8) {test fallito};
\draw[freccia] (e) -- (o);
\draw[freccia] (o) -- (q);
\draw[freccia] (q) -- node[right, font=\footnotesize] {sì} (ok);
\draw[freccia] (q) -- node[right, font=\footnotesize] {no} (ko);
\end{tikzpicture}
\caption{Un test. Si esegue il programma su dati di prova, si osserva il comportamento e lo si confronta con quello
atteso.}
\label{fig:test-ciclo}
\end{figure}

I tipi di testing sono quattro.
\begin{description}[leftmargin=0pt, style=unboxed]
  \item[Functional testing.] Si collauda la funzionalità del sistema complessivo per scoprire i bug (\rif{sec:test-funzionale}).
  \item[User testing.] Si verifica che il prodotto sia utile e usabile per gli utenti finali.
    \begin{itemize}
      \item \termine{Alpha testing}: si chiede agli utenti ``volete davvero queste feature?''.
      \item \termine{Beta testing}: si distribuisce agli utenti una versione preliminare del prodotto, per controllarne
        l'usabilità e la capacità di funzionare insieme ad altri sistemi (\ingl{interoperability}).
    \end{itemize}
  \item[Performance e load testing.] Si verifica che il sistema risponda in fretta alle richieste. Si verifica anche
    che regga carichi diversi e che \emph{scali} senza problemi quando il carico cresce.
  \item[Security testing.] Si cercano le vulnerabilità che un attaccante può sfruttare (\rif{sec:test-sicurezza}).
\end{description}

\begin{attenzione}[Il testing non è una verifica]
\uline{Il testing può mostrare la presenza di bug, mai la loro assenza} (Edsger W. Dijkstra). \emph{Software
testing} $\neq$ \emph{software verification}: lo mostra \figrif{fig:test-verifica}.
\end{attenzione}

\begin{figure}[ht]
\centering
\begin{tikzpicture}[scale=0.9]
\foreach \i in {0,...,7} {
  \pgfmathsetmacro{\x}{0.7*\i}
  \pgfmathsetmacro{\y}{-0.55*\i}
  \ifnum\i=1
    \fill[rosso] (\x,\y) rectangle (\x+0.85,\y-0.2);
  \else
    \fill[inchiostro] (\x,\y) rectangle (\x+0.85,\y-0.2);
  \fi
}
\foreach \i in {0,2,4,6} {
  \pgfmathsetmacro{\x}{0.7*\i}
  \pgfmathsetmacro{\y}{-0.55*\i-0.1}
  \draw[->, verde] (\x-1.1,\y) -- (\x-0.05,\y);
  \node[font=\footnotesize, text=verde, anchor=east] at (\x-1.1,\y) {test};
}
\node[font=\Large, text=rosso] at (7.2,-1.9) {$\neq$};
\node[draw, text width=3.6cm, align=center, font=\small] at (10.0,-1.9) {dimostrazione:\\vale per \emph{tutti}\\gli input};
\end{tikzpicture}
\caption{Testing e verifica. A sinistra ogni gradino è un possibile input; i test (frecce verdi) ne controllano solo
alcuni. Il gradino rosso, che contiene un bug, non è toccato da nessun test e il bug resta nascosto. A destra una
dimostrazione, come quella di un teorema, vale per tutti i casi.}
\label{fig:test-verifica}
\end{figure}

\subsubsection{Perché il testing non dimostra la correttezza\aggiunta}
Un programma ha di solito un numero enorme di input possibili, e i test ne provano solo alcuni. Se tutti i test
riescono, si sa che il programma funziona \emph{su quegli input}. Sugli altri non si sa nulla. La
\termine{verifica} del software invece dimostra con un ragionamento formale che il programma è corretto per
\emph{ogni} input, come si dimostra un teorema.

\section{Functional testing}
\label{sec:test-funzionale}

Il functional testing segue alcune regole generali.
\begin{itemize}
  \item Si prepara un ampio insieme di test, in modo che \textbf{tutto il codice sia eseguito almeno una volta}.
  \item Il testing comincia il giorno in cui si comincia a scrivere il codice.
  \item I test automatici (\rif{sec:test-automazione}) rendono più semplice il ciclo sviluppo-test.
\end{itemize}

Il functional testing è un'attività \emph{a fasi} (\figrif{fig:test-fasi}): unit testing, feature testing, system
testing, release testing. Il ciclo si ripete a ogni nuova versione.

\begin{figure}[ht]
\centering
\begin{tikzpicture}
\node[blocco] (u) at (0,1.5) {Unit testing};
\node[blocco] (f) at (2.6,0) {Feature testing};
\node[blocco] (s) at (0,-1.5) {System testing};
\node[blocco] (r) at (-2.6,0) {Release testing};
\node[font=\footnotesize] (st) at (0,2.7) {inizio};
\draw[freccia] (st) -- (u);
\draw[freccia] (u) to[bend left=20] (f);
\draw[freccia] (f) to[bend left=20] (s);
\draw[freccia] (s) to[bend left=20] (r);
\draw[freccia] (r) to[bend left=20] (u);
\end{tikzpicture}
\caption{Le fasi del functional testing: unità, feature, sistema, rilascio. Dopo il rilascio il ciclo riparte con le
modifiche successive.}
\label{fig:test-fasi}
\end{figure}

\subsection{Unit testing}
\label{sec:test-unit}

Lo \termine{unit testing} collauda le unità del programma (per esempio una funzione o un metodo) \emph{in
isolamento}. Si basa su un principio.

\begin{idea}[Principio dello unit testing]
Si prenda un insieme di input che hanno alcune caratteristiche in comune. Se un'unità di programma si comporta come
previsto su quell'insieme, si comporterà allo stesso modo su un insieme più grande con le stesse caratteristiche.
\end{idea}

Per esempio, se il programma si comporta correttamente sull'insieme di input $\{1, 5, 17, 45, 99\}$, si può concludere
che tratterà correttamente tutti gli altri interi tra 1 e 99.

Da qui nasce il metodo delle \termine{equivalence partition}.
\begin{itemize}
  \item Si individuano le \emph{partizioni di equivalenza}: insiemi di input che il codice tratta allo stesso modo.
  \item Si collauda il programma con più input presi da ogni partizione.
  \item I programmatori sbagliano spesso ai confini: si individuano i \textbf{confini delle partizioni} e si scelgono
    input proprio su quei confini.
\end{itemize}

\begin{esempio}[Partizioni e confini\aggiunta]
Una funzione accetta gli interi da 1 a 99. Le partizioni sono tre: gli interi minori di 1 (da rifiutare), quelli da 1
a 99 (da accettare) e quelli maggiori di 99 (da rifiutare). Si sceglie qualche valore dentro ogni partizione, per
esempio $-5$, $45$ e $150$. Poi si aggiungono i valori ai confini: $0$ e $1$, $99$ e $100$. Un errore comune come
scrivere \texttt{x < 99} invece di \texttt{x <= 99} viene scoperto solo dal test con 99.
\end{esempio}

\tabrif{tab:test-linee-guida} riporta le linee guida per scegliere i casi di test.

\begin{table}[ht]
\centering
\small
\caption{Linee guida per lo unit testing.}
\label{tab:test-linee-guida}
\begin{tabularx}{\linewidth}{>{\raggedright\arraybackslash}p{3.4cm}X}
\toprule
Linea guida & Spiegazione \\
\midrule
Test edge cases & Se la partizione ha limiti inferiore e superiore (lunghezza di stringhe, numeri, ecc.), scegliere
  input ai bordi dell'intervallo. \\
Force errors & Scegliere input che costringano il sistema a generare tutti i messaggi d'errore. Scegliere input che
  dovrebbero produrre output non validi. \\
Fill buffers & Scegliere input che facciano traboccare tutti i buffer di input. \\
Repeat yourself & Ripetere più volte lo stesso input o la stessa sequenza di input. \\
Overflow and underflow & Se il programma fa calcoli numerici, scegliere input che lo portino a calcolare numeri molto
  grandi o molto piccoli. \\
Don't forget null and zero & Con puntatori o stringhe, provare sempre puntatori nulli e stringhe nulle. Con le
  sequenze, provare la sequenza vuota. Con input numerici, provare sempre lo zero. \\
Keep count & Con liste e trasformazioni di liste, tenere il conto degli elementi di ogni lista e controllare che i
  conteggi restino coerenti dopo ogni trasformazione. \\
One is different & Con le sequenze, provare sempre sequenze di un solo elemento. \\
\bottomrule
\end{tabularx}
\end{table}

\subsection{Feature testing}
\label{sec:test-feature}

Una feature del prodotto realizza una funzionalità utile all'utente (\rif{sec:feature-intro}). Va collaudata per
mostrare che la funzionalità:
\begin{enumerate}
  \item è realizzata come previsto;
  \item soddisfa i bisogni reali degli utenti.
\end{enumerate}
Una feature è di solito realizzata da più unità di programma che interagiscono. Per questo i test sono di due tipi.
\begin{description}[leftmargin=0pt, style=unboxed]
  \item[Interaction test.] Collaudano le interazioni tra unità, spesso sviluppate da sviluppatori diversi. Possono
    rivelare anche bug nelle singole unità che lo unit testing non aveva fatto emergere.
  \item[Usefulness test.] Verificano che la feature realizzi ciò che gli utenti probabilmente vogliono. Il product
    manager dovrebbe partecipare alla loro progettazione.
\end{description}

I feature test si ricavano dalle user story (\rif{sec:feature-storie}). \tabrif{tab:test-storie} mostra un esempio.

\begin{table}[ht]
\centering
\small
\caption{Dalle user story ai feature test. In alto tre user story sulla registrazione; in basso i test che se ne
ricavano.}
\label{tab:test-storie}
\begin{tabularx}{\linewidth}{>{\raggedright\arraybackslash}p{3.2cm}X}
\toprule
Storia & User story \\
\midrule
User registration & As a user, I want to be able to log in without creating a new account so that I don't have to
  remember another login ID and password. \\
Information sharing & As a user, I want to know what information you will share with other companies. I want to be
  able to cancel my registration if I don't want to share this information. \\
Email choice & As a user, I want to be able to choose the types of email that I'll get from you when I register for
  an account. \\
\midrule
Test & Descrizione \\
\midrule
Initial login screen & Il clic sul link ``Sign-in with Google'' mostra correttamente la richiesta delle credenziali
  Google. Il login si completa se l'utente ha già fatto l'accesso a Google. \\
Incorrect credentials & Se l'utente inserisce credenziali Google sbagliate, compaiono il messaggio d'errore e la
  schermata per riprovare. \\
Shared information & Compaiono le informazioni condivise con Google, con le opzioni annulla e conferma. Se si sceglie
  annulla, la registrazione è annullata. \\
Email opt-in & All'utente si offre un menu di opzioni sulle email e può sceglierne più di una. Se non sceglie
  nessuna opzione, non è iscritto a nessuna email. \\
\bottomrule
\end{tabularx}
\end{table}

\subsection{System testing}
\label{sec:test-sistema}

Il \termine{system testing} collauda il sistema nel suo insieme, per quattro scopi:
\begin{itemize}
  \item scoprire interazioni inattese o indesiderate tra le feature;
  \item scoprire se le feature lavorano bene insieme per sostenere ciò che gli utenti vogliono davvero fare;
  \item assicurarsi che il sistema funzioni come previsto nei diversi ambienti in cui sarà usato;
  \item collaudare reattività, throughput, sicurezza e gli altri attributi di qualità.
\end{itemize}
Un consiglio: usare scenari e user story (capitolo~\ref{cap:feature}) per individuare i percorsi completi
(\ingl{end-to-end}) degli utenti nel sistema.

\subsection{Release testing}
\label{sec:test-release}

Il \termine{release testing} collauda un sistema destinato al rilascio ai clienti.
\begin{itemize}
  \item Collauda il sistema nel suo \textbf{ambiente operativo reale}, invece che in un ambiente di test.
  \item Lo scopo è decidere se il sistema è abbastanza buono da essere rilasciato, \emph{non} trovare bug.
\end{itemize}
Serve perché preparare il rilascio significa impacchettare il sistema per la distribuzione, installare il software e
le librerie che usa e configurare i parametri. In ognuno di questi passi si possono commettere errori. Se si
distribuisce sul cloud si può usare un processo di rilascio automatico e continuo (capitolo~\ref{cap:devops}).

\section{Test automation}
\label{sec:test-automazione}

Nelle aziende che sviluppano prodotti si usa molto la test automation, già incontrata tra le pratiche di Scrum
(\rif{sec:agile-sprint}).
\begin{itemize}
  \item Esistono framework di test per tutti i linguaggi di programmazione più usati.
  \item Un test automatico è un programma eseguibile che controlla che il software restituisca il risultato atteso
    per certi dati di input.
\end{itemize}

\begin{algoritmo}[Arrange, Action, Assert]
Una buona pratica divide ogni test automatico in tre parti.
\begin{enumerate}
  \item \textbf{Arrange}: si prepara il sistema per eseguire il test. Si definiscono i parametri del test e, se
    servono, i \emph{mock object}, oggetti finti che sostituiscono le parti del sistema non collaudate.
  \item \textbf{Action}: si chiama l'unità da collaudare con i parametri del test.
  \item \textbf{Assert}: si afferma ciò che deve valere se il test è eseguito con successo.
\end{enumerate}
\end{algoritmo}

\begin{esempio}[Un test per il calcolo degli interessi]
\begin{verbatim}
#Arrange - set up the test parameters
p = 0
r = 3
n = 31
result_should_be = 0

#Action - Call the method to be tested
interest = interest_calculator (p, r, n)

#Assert - test what should be true
self.assertEqual (result_should_be, interest)
\end{verbatim}
Con un capitale \texttt{p} uguale a zero l'interesse deve essere zero, qualunque siano il tasso \texttt{r} e il numero
di giorni \texttt{n}. È anche un caso della linea guida ``don't forget null and zero'' (\tabrif{tab:test-linee-guida}).
\end{esempio}

\begin{attenzione}[Anche il codice dei test può avere bug]
Per ridurre il rischio di errori nei test:
\begin{enumerate}
  \item si scrivono test il più semplici possibile;
  \item si rivedono tutti i test insieme al codice che collaudano.
\end{enumerate}
\end{attenzione}

I tipi di test non si automatizzano con la stessa facilità (\figrif{fig:test-piramide}).
\begin{itemize}
  \item Gli \textbf{unit test} sono i più facili da automatizzare.
  \item Buoni unit test riducono, ma non eliminano, il bisogno di feature test.
  \item I \textbf{feature test} basati sull'interfaccia grafica (GUI) sono costosi da automatizzare: è preferibile
    collaudare attraverso le API. Servono più asserzioni per controllare che la feature sia stata eseguita come
    previsto.
  \item Il \textbf{system testing} collauda il sistema facendo le veci di un utente: esegue sequenze di azioni utente.
    Il system testing manuale è noioso e soggetto a errori. Esistono strumenti che registrano una serie di azioni e la
    ripetono in automatico.
\end{itemize}

\begin{figure}[ht]
\centering
\begin{tikzpicture}
% base larga 7, altezza 5,5: semilarghezza a quota y = 3,5 (1 - y/5,5)
\draw (0,0) -- (7,0) -- (3.5,5.5) -- cycle;
\draw (3.5-2.164,2.1) -- (3.5+2.164,2.1);
\draw (3.5-1.145,3.7) -- (3.5+1.145,3.7);
\node[font=\small] at (3.5,1.0) {Unit tests};
\node[font=\small] at (3.5,2.9) {Feature tests};
\node[font=\footnotesize, align=center] at (3.5,4.25) {System\\tests};
\node[font=\small, anchor=west] at (6.5,1.0) {70\%};
\node[font=\small, anchor=west] at (5.3,2.9) {20\%};
\node[font=\small, anchor=west] at (4.5,4.25) {10\%};
\draw[freccia, blu] (-0.6,5.2) -- (-0.6,0.2);
\node[font=\footnotesize, text=blu, align=right, anchor=east] at (-0.8,2.7) {più automazione,\\costi minori};
\end{tikzpicture}
\caption{La piramide dei test. Alla base gli unit test, che sono il 70\% dei test; poi i feature test (20\%) e in cima
i system test (10\%). Scendendo verso la base i test sono più facili da automatizzare e costano meno.}
\label{fig:test-piramide}
\end{figure}

\section{Test-driven development}
\label{sec:test-tdd}

Il \termine{test-driven development} (TDD) è una pratica di Extreme Programming (\rif{sec:agile-xp}): \textbf{prima si
scrive il test eseguibile, poi il codice}. Si scrivono piccoli unit test (\ingl{mini-unit test}) per una parte della
funzionalità alla volta (\figrif{fig:test-tdd}).

\begin{algoritmo}[Ciclo del TDD]
\begin{enumerate}
  \item Individuare una nuova funzionalità.
  \item Individuare una sua realizzazione parziale, abbastanza piccola da scriverne i test.
  \item Scrivere i mini-unit test per quella parte.
  \item Scrivere uno \emph{stub} del codice, un abbozzo che farà fallire il test.
  \item Eseguire tutti i test automatici: il nuovo test fallisce.
  \item Scrivere il codice che dovrebbe far passare il test che fallisce.
  \item Eseguire tutti i test automatici. Se qualcuno fallisce, tornare al passo 6.
  \item Quando tutti i test passano, rifattorizzare il codice se serve.
  \item Se la funzionalità è incompleta, tornare al passo 2; se è completa, tornare al passo 1.
\end{enumerate}
\end{algoritmo}

\begin{figure}[ht]
\centering
\begin{tikzpicture}[b/.style={blocco, font=\footnotesize, text width=2.6cm}]
\node[b] (a) at (0,4.2) {Individuare una nuova funzionalità};
\node[b] (b) at (0,2.4) {Individuare una realizzazione parziale};
\node[b] (c) at (4.6,2.4) {Scrivere i mini-unit test};
\node[b] (d) at (9.2,2.4) {Scrivere uno stub che farà fallire il test};
\node[b] (e) at (9.2,0.6) {Eseguire tutti i test automatici};
\node[b] (f) at (9.2,-1.2) {Scrivere il codice che fa passare il test};
\node[b] (g) at (4.6,-1.2) {Eseguire tutti i test automatici};
\node[b] (h) at (0,-1.2) {Rifattorizzare il codice se serve};
\node[font=\footnotesize] (s) at (0,5.5) {inizio};
\draw[freccia] (s) -- (a);
\draw[freccia] (a) -- (b);
\draw[freccia] (b) -- (c);
\draw[freccia] (c) -- (d);
\draw[freccia] (d) -- (e);
\draw[freccia] (e) -- (f);
\draw[freccia] ([yshift=-4pt]f.west) -- ([yshift=-4pt]g.east);
\draw[freccia] ([yshift=4pt]g.east) -- node[above, font=\scriptsize] {test fallito} ([yshift=4pt]f.west);
\draw[freccia] (g) -- node[below, font=\scriptsize] {tutti passano} (h);
\draw[freccia] (h) -- node[right, font=\scriptsize, align=left] {funzionalità\\incompleta} (b);
\draw[freccia] (h.west) -- ++(-0.5,0) |- node[pos=0.25, left, font=\scriptsize, align=right] {funzionalità\\completa} (a.west);
\end{tikzpicture}
\caption{Il ciclo del test-driven development. Si scrivono prima i test per una piccola parte della funzionalità,
poi il codice che li fa passare. Se un test fallisce si torna al codice; quando tutti passano si rifattorizza e si
passa alla parte successiva, o alla funzionalità successiva.}
\label{fig:test-tdd}
\end{figure}

\begin{table}[ht]
\centering
\small
\caption{Vantaggi e svantaggi del test-driven development.}
\label{tab:test-tdd}
\begin{tabularx}{\linewidth}{>{\raggedright\arraybackslash}X>{\raggedright\arraybackslash}X}
\toprule
Vantaggi & Svantaggi \\
\midrule
Approccio sistematico: i test sono legati chiaramente a sezioni di codice e non restano sezioni non collaudate. &
  È difficile applicare il TDD al system testing. \\
I test aiutano a capire il codice del programma. & Il TDD scoraggia i cambiamenti radicali del programma. \\
Il debugging è più semplice e incrementale. & Il TDD porta a concentrarsi sui test invece che sul problema da
  risolvere. \\
Il codice è (forse) più semplice. & Il TDD porta a pensare troppo ai dettagli di implementazione invece che alla
  struttura complessiva del programma. \\
\bottomrule
\end{tabularx}
\end{table}

\section{Security testing}
\label{sec:test-sicurezza}

Il \termine{security testing} ha due obiettivi:
\begin{itemize}
  \item trovare le vulnerabilità che un attaccante può sfruttare;
  \item fornire prove convincenti che il sistema è abbastanza sicuro.
\end{itemize}

Trovare vulnerabilità è più difficile che trovare bug, per tre motivi.
\begin{itemize}
  \item Bisogna collaudare qualcosa che il software \emph{non} deve fare: il numero di test possibili è potenzialmente
    infinito.
  \item I normali test funzionali possono non rivelare le vulnerabilità.
  \item Lo stack software su cui si basa il prodotto (sistema operativo, librerie, database, ecc.) può contenere
    vulnerabilità.
\end{itemize}

Un security testing completo richiede conoscenze specialistiche. Per esempio si coinvolgono specialisti esterni per il
\termine{penetration testing}, cioè per tentare attacchi reali al sistema; è costoso. Si adotta quindi un
\textbf{approccio basato sul rischio}:
\begin{itemize}
  \item si individuano i principali rischi di sicurezza del prodotto;
  \item si sviluppano test che dimostrino che il prodotto si protegge da quei rischi.
\end{itemize}
Alcuni di questi test si possono automatizzare. Altri richiedono inevitabilmente di controllare a mano comportamenti o
file.

\begin{esempio}[Rischi di sicurezza]
\begin{itemize}
  \item Un attaccante non autorizzato accede al sistema usando credenziali autorizzate.
  \item Una persona autorizzata accede a risorse che le sono vietate.
  \item Il sistema di autenticazione non rileva un attaccante non autorizzato.
  \item L'attaccante accede al database con un attacco di SQL poisoning (\rif{sec:sic-injection}).
  \item Le sessioni HTTP sono gestite in modo scorretto.
  \item I cookie di sessione HTTP sono rivelati a un attaccante (\rif{sec:sic-sessione}).
  \item Dati riservati non sono cifrati.
  \item Le chiavi di cifratura finiscono in mano a possibili attaccanti.
\end{itemize}
Gli attacchi citati sono descritti nel capitolo~\ref{cap:sicurezza}.
\end{esempio}

\begin{idea}[Pensare come un attaccante]
Nel security testing bisogna pensare come un attaccante, non come un normale utente finale:
\begin{itemize}
  \item provare di proposito a fare la cosa sbagliata;
  \item ripetere più volte le stesse azioni.
\end{itemize}
\end{idea}

\section{Limiti del testing e code review}
\label{sec:test-review}

Il testing ha tre limiti.
\begin{enumerate}
  \item Si collauda il codice rispetto a ciò che \emph{si è capito} che il codice deve fare. Se se ne è frainteso lo
    scopo, l'errore si ritrova sia nel codice sia nei test.
  \item Il testing può non coprire tutto il codice scritto. Il TDD sposta il problema: invece di codice non collaudato,
    si rischia codice incompleto.
  \item Il testing non dice nulla sugli altri attributi di un programma, come leggibilità, struttura ed evolvibilità.
\end{enumerate}

Per questo le \termine{code review} \textbf{completano il testing}. Una code review procede in quattro fasi
(\figrif{fig:test-review}).

\begin{figure}[ht]
\centering
\begin{tikzpicture}[b/.style={blocco, font=\footnotesize, minimum width=2.1cm, text width=1.9cm},
  g/.style={draw, fill=grigiochiaro, inner sep=5pt}]
% fase 1
\node[b] (a1) at (0,1.2) {Set up review};
\node[b] (a2) at (0,0) {Prepare code};
\node[b] (a3) at (0,-1.2) {Distribute code/tests};
% fase 2
\node[b] (b1) at (3.4,0.6) {Check code};
\node[b] (b2) at (3.4,-0.6) {Write review report};
% fase 3
\node[draw, ellipse, font=\footnotesize, minimum width=2.3cm, minimum height=1.0cm] (c) at (6.9,0) {Discussion};
% fase 4
\node[b] (d1) at (10.3,0.6) {Prepare to-do list};
\node[b] (d2) at (10.3,-0.6) {Make code changes};
\begin{scope}[on background layer]
\node[g, fit=(a1)(a3)] (fa) {};
\node[g, fit=(b1)(b2)] (fb) {};
\node[g, fit=(d1)(d2)] (fd) {};
\end{scope}
\draw[freccia] (fa) -- (fb);
\draw[freccia] (fb) -- (c);
\draw[freccia] (c) -- (fd);
\node[font=\footnotesize] at (0,2.2) {Programmer};
\node[font=\footnotesize] at (3.4,1.6) {Reviewer};
\node[font=\footnotesize, align=center] at (6.9,1.0) {Reviewer e\\programmer};
\node[font=\footnotesize] at (10.3,1.6) {Programmer};
\node[font=\footnotesize] at (0,-2.2) {Review preparation};
\node[font=\footnotesize] at (3.4,-1.6) {Code checking};
\node[font=\footnotesize] at (6.9,-1.6) {Review};
\node[font=\footnotesize] at (10.3,-1.6) {Follow-up};
\end{tikzpicture}
\caption{Le quattro fasi di una code review. Il programmatore prepara la review e distribuisce codice e test; il
revisore controlla il codice e scrive un rapporto; i due ne discutono; infine il programmatore prepara la lista delle
cose da fare e modifica il codice. In alto chi agisce, in basso il nome della fase.}
\label{fig:test-review}
\end{figure}

\begin{itemize}
  \item Spesso c'è un solo revisore, del team DevOps (capitolo~\ref{cap:devops}) oppure no.
  \item Il revisore può commentare anche la leggibilità e la comprensibilità del codice.
  \item Una sessione di review dovrebbe concentrarsi su 200--400 righe di codice.
  \item Le review possono partire automaticamente dai commit su un repository condiviso.
  \item Esistono molti strumenti per le code review, spesso usati insieme a sistemi di messaggistica come Slack.
\end{itemize}

Si usa spesso una \emph{checklist}. Per esempio:
\begin{itemize}
  \item si usano nomi significativi per variabili e funzioni? (generale)
  \item sono stati considerati tutti gli errori sui dati, con i relativi test? (generale)
  \item tutte le eccezioni sono gestite esplicitamente? (generale)
  \item si usano parametri di default delle funzioni? (Python)
  \item i tipi sono usati in modo coerente? (Python)
  \item il livello di indentazione è corretto? (Python)
\end{itemize}

\begin{esempio}[Code review a Spotify]
\label{es:test-spotify}
A Spotify il lavoro è organizzato in \emph{squad}, piccoli team autonomi (\rif{sec:devops-spotify}).
\begin{itemize}
  \item Un \termine{chapter} riunisce i membri dei team che lavorano in un'area specifica: sviluppatori di front office,
    sviluppatori di back office, amministratori di database, tester.
  \item Una \termine{guild} è una comunità di persone di tutta l'organizzazione con interessi comuni (tecnologie web,
    test automation), che vogliono condividere conoscenze, strumenti, codice e pratiche.
\end{itemize}
I chapter (a volte le guild) fanno le code review per le squad. Per fare il merge servono due ``+1'', cioè due
approvazioni.
\end{esempio}

\begin{puntichiave}
\begin{itemize}
  \item Testing: eseguire il programma su dati che simulano gli input e confrontare il comportamento con quello atteso.
    Mostra la presenza di bug, mai la loro assenza: testing $\neq$ verification.
  \item Tipi: functional, user (alpha, beta), performance e load, security.
  \item Functional testing a fasi:
    \begin{itemize}
      \item unit: in isolamento, con equivalence partition e confini;
      \item feature: interaction e usefulness test, ricavati dalle user story;
      \item system: il sistema intero, percorsi end-to-end;
      \item release: nell'ambiente reale, per decidere se rilasciare, non per trovare bug.
    \end{itemize}
  \item Test automatici in tre parti: arrange, action, assert. Piramide: 70\% unit, 20\% feature, 10\% system.
  \item TDD: prima il test, poi il codice, una piccola parte alla volta. Vantaggi e svantaggi in \tabrif{tab:test-tdd}.
  \item Security testing: basato sul rischio; pensare come un attaccante.
  \item Il testing non vede gli errori di comprensione né la qualità del codice: le code review lo completano.
\end{itemize}
\end{puntichiave}

\begin{domande}
\begin{enumerate}
  \item Che cos'è il testing? Perché il testing non è una verifica?
  \item Quali sono i tipi di testing? Che differenza c'è tra alpha e beta testing?
  \item Enunciare il principio dello unit testing. Che cosa sono le equivalence partition e perché si testano i confini?
  \item Che differenza c'è tra interaction test e usefulness test? Come si ricavano i feature test dalle user story?
  \item Quali sono gli scopi del system testing? In che cosa il release testing è diverso?
  \item Come si struttura un test automatico? Perché la piramide dei test ha quella forma?
  \item Descrivere il ciclo del TDD. Quali sono i suoi vantaggi e svantaggi?
  \item Perché trovare vulnerabilità è più difficile che trovare bug? Che cos'è l'approccio basato sul rischio?
  \item Quali sono i limiti del testing? Come si svolge una code review?
\end{enumerate}
\end{domande}
